% ============================================================
\chapter{Vector Spaces and Subspaces}
\label{ch:vector-spaces-subspaces}
% ============================================================

\section*{Chapter Preview}
\addcontentsline{toc}{section}{Chapter Preview}

This chapter gives the formal definition that Chapter~\ref{ch:vectors-data}
postponed. A vector space is a set with two operations: addition and scalar
multiplication. A subspace is a smaller vector space sitting inside a larger
one.

\section{Vector Spaces}

\begin{definition}
  A \emph{vector space} $V$ over a field $\F$ is a set equipped with
  \begin{itemize}
    \item addition $+:V\times V\to V$,
    \item scalar multiplication $\F\times V\to V$,
  \end{itemize}
  satisfying the following axioms for all $u,v,w\in V$ and all
  $\lambda,\mu\in\F$:
  \begin{enumerate}
    \item $(u+v)+w=u+(v+w)$.
    \item $u+v=v+u$.
    \item There is a vector $0\in V$ such that $v+0=v$.
    \item For every $v\in V$ there is a vector $-v\in V$ such that
      $v+(-v)=0$.
    \item $(\lambda\mu)v=\lambda(\mu v)$.
    \item $(\lambda+\mu)v=\lambda v+\mu v$.
    \item $\lambda(u+v)=\lambda u+\lambda v$.
    \item $1v=v$.
  \end{enumerate}
\end{definition}

The definition is long, but most examples are familiar. In $\R^n$, addition
and scalar multiplication are componentwise. In a polynomial space, they are
coefficientwise. In a function space, they are pointwise.

\begin{example}[Polynomials of degree at most $n$]
Let
\[
  \mathcal P_n
  =
  \{a_0+a_1x+\cdots+a_nx^n : a_i\in\R\}.
\]
Adding polynomials and multiplying them by real scalars stays inside
$\mathcal P_n$. The zero vector is the zero polynomial.
\end{example}

\section{Consequences of the Axioms}

The axioms imply facts that we usually take for granted.

\begin{proposition}
  The zero vector in a vector space is unique.
\end{proposition}

\begin{proof}
  Suppose $0$ and $0'$ are both zero vectors. Since $0'$ is a zero vector,
  $0+0'=0$. Since $0$ is a zero vector, $0+0'=0'$. Therefore $0=0'$.
\end{proof}

\begin{proposition}
  For every vector $v$, one has $0_{\F}v=0_V$.
\end{proposition}

\begin{proof}
  Since $0_{\F}+0_{\F}=0_{\F}$,
  \[
    0_{\F}v=(0_{\F}+0_{\F})v=0_{\F}v+0_{\F}v.
  \]
  Add the additive inverse of $0_{\F}v$ to both sides.
\end{proof}

\section{Subspaces}

\begin{definition}
  Let $V$ be a vector space over $\F$. A subset $W\subseteq V$ is a
  \emph{subspace} if $W$ is itself a vector space using the addition and scalar
  multiplication inherited from $V$.
\end{definition}

There is a simple test.

\begin{theorem}[Subspace test]
  A nonempty subset $W\subseteq V$ is a subspace if and only if for all
  $u,v\in W$ and all $\lambda,\mu\in\F$,
  \[
    \lambda u+\mu v\in W.
  \]
\end{theorem}

\begin{proof}
  If $W$ is a subspace, it is closed under scalar multiplication and addition,
  so $\lambda u+\mu v\in W$.

  Conversely, assume the displayed condition. Taking $\lambda=\mu=1$ gives
  closure under addition. Taking $\mu=0$ gives closure under scalar
  multiplication. Since $W$ is nonempty, choose $w\in W$; then
  $0w=0\in W$. Additive inverses follow from $-w=(-1)w\in W$. The remaining
  vector-space axioms are inherited from $V$.
\end{proof}

\begin{example}[A plane through the origin]
The set
\[
  W=\{(x,y,z)\in\R^3:x+y=0\}
\]
is a subspace of $\R^3$. If $u$ and $v$ satisfy the equation, then any linear
combination $\lambda u+\mu v$ also satisfies it.
\end{example}

\begin{example}[An affine plane is not a subspace]
The set
\[
  A=\{(x,y,z)\in\R^3:x+y=1\}
\]
is not a subspace. It does not contain the zero vector.
\end{example}

\section{Python: Testing Closure}

\begin{lstlisting}
import numpy as np

def in_plane(v, tol=1e-10):
    return abs(v[0] + v[1]) < tol

def sample_plane():
    x = np.random.randn()
    z = np.random.randn()
    return np.array([x, -x, z])

ok = True
for _ in range(200):
    u = sample_plane()
    v = sample_plane()
    a, b = np.random.randn(), np.random.randn()
    ok = ok and in_plane(a*u + b*v)

print("closure test passed:", ok)
\end{lstlisting}

\begin{warningbox}
A random numerical test can find counterexamples, but it is not a proof. The
proof is the symbolic argument that arbitrary linear combinations stay inside
the set.
\end{warningbox}

\section*{Chapter Summary}
\addcontentsline{toc}{section}{Chapter Summary}

\begin{itemize}
  \item A vector space is a set with addition and scalar multiplication
    satisfying eight axioms.
  \item The zero vector is unique, and $0_{\F}v=0_V$.
  \item A subspace is a subset that is closed under linear combinations.
  \item Passing random numerical tests is not the same as proof.
\end{itemize}

\section*{Where This Appears Later}
\addcontentsline{toc}{section}{Where This Appears Later}

Subspaces reappear as spans, kernels, images, row spaces, column spaces,
eigenspaces, and low-rank approximation spaces.

\begin{labbox}
Lab 1 includes a stochastic axiom checker and subspace tests. Use it to see
how numerical experiments can support, but not replace, proof.
\end{labbox}

\section*{Exercises}
\addcontentsline{toc}{section}{Exercises}

\subsection*{A. Theory}

\begin{exercise}
\exerciseid{3.A1}
Prove that $\R^n$ with componentwise addition and scalar multiplication is a
vector space. You may group the axioms by saying which properties follow from
the corresponding properties of real numbers.
\end{exercise}

\begin{exercise}
\exerciseid{3.A2}
Prove that the zero vector in a vector space is unique.
\end{exercise}

\begin{exercise}
\exerciseid{3.A3}
Let $A$ be an $m\times n$ matrix. Prove that
\[
  \{x\in\R^n:Ax=0\}
\]
is a subspace of $\R^n$.
\end{exercise}

\begin{exercise}
\exerciseid{3.A4}
Find the false step in this argument: ``The set
$\{(x,y):x+y=1\}$ is a subspace because the equation is linear.''
\end{exercise}

\subsection*{B. By Hand}

\begin{exercise}
\exerciseid{3.B1}
Decide whether each subset of $\R^2$ is a subspace:
\[
  \{(x,y):y=2x\},\quad
  \{(x,y):y=2x+1\},\quad
  \{(x,y):xy=0\}.
\]
Give a short reason for each answer.
\end{exercise}

\begin{exercise}
\exerciseid{3.B2}
Let
\[
  W=\{(x,y,z)\in\R^3:x-y+z=0\}.
\]
Find three nonzero vectors in $W$. Then check by hand that the sum of any two
of your vectors is still in $W$.
\end{exercise}

\begin{exercise}
\exerciseid{3.B3}
Give a vector in the set $\{(x,y,z):x+y+z=1\}$ and show directly why doubling
it leaves the set.
\end{exercise}

\subsection*{C. Python / NumPy}

\begin{exercise}
\exerciseid{3.C1}
Write a function \texttt{in\_plane(v)} that tests whether
$v_1+v_2+v_3=0$ up to tolerance. Generate random examples in the plane and
test closure under random linear combinations.
\end{exercise}

\begin{exercise}
\exerciseid{3.C2}
Write a stochastic test for the unit circle in $\R^2$. Use it to find evidence
that the unit circle is not a subspace.
\end{exercise}

\subsection*{D. Challenge}

\begin{exercise}
\exerciseid{3.D1}
Let $W_1$ and $W_2$ be subspaces of $V$. Prove that $W_1\cap W_2$ is a
subspace. Is $W_1\cup W_2$ always a subspace?
\end{exercise}

